\documentclass[11pt]{article}

\usepackage[margin=1.15in]{geometry}
\usepackage{amsmath,amssymb,amsthm,mathtools}
\usepackage[T1]{fontenc}
\usepackage{lmodern}
\usepackage{microtype}

\newtheorem{theorem}{Theorem}
\newtheorem{lemma}[theorem]{Lemma}
\newtheorem{proposition}[theorem]{Proposition}
\newtheorem{corollary}[theorem]{Corollary}
\newtheorem{remark}[theorem]{Remark}

\newcommand{\F}{\mathbb F}
\newcommand{\E}{\mathbb E}
\newcommand{\Prb}{\mathbb P}
\newcommand{\Inf}{\operatorname{Inf}}
\newcommand{\AS}{\operatorname{AS}}
\newcommand{\sens}{\operatorname{S}}

\title{A monotone Boolean construction with average--max sensitivity exponent \(2/3\)}
\author{Slava Naprienko}
\date{}

\begin{document}
\maketitle

\begin{abstract}
The Servedio--Tan open problem asks whether every monotone Boolean function \(f\) satisfies \(\AS(f)=o(\sens(f))\) \cite{FHH+14,ScTa13}. We construct a sequence of monotone Boolean functions \(H_k\) with \(\sens(H_k)\to\infty\) and
\[
    \liminf_{k\to\infty}\frac{\log\AS(H_k)}{\log\sens(H_k)}\ge \frac{2}{3}.
\]
This gives an explicit monotone family with exponent \(2/3\) for the Servedio--Tan average-vs-max sensitivity problem, improving the previously recorded lower-bound exponent
\(
    \alpha_{\mathrm{OS}}=\log_2(6-2\sqrt 5)\approx 0.6115
\)
of O'Donnell and Servedio \cite{OS07}.
\end{abstract}

\section{Introduction}

Let \(f:\{0,1\}^n\to\{0,1\}\) be a monotone Boolean function. Two standard complexity measures of \(f\) are its average sensitivity (also called total influence) \(\AS(f)=\Inf_{1/2}(f)\), the average over a uniform input of the number of pivotal coordinates, and its (maximum) sensitivity \(\sens(f)\), the corresponding worst-case quantity. Trivially \(\AS(f)\le\sens(f)\), and Servedio and Tan (recorded under the heading \emph{Average versus Max Sensitivity for Monotone Functions} in \cite{FHH+14}) asked whether monotone functions in fact satisfy
\[
    \AS(f)=o(\sens(f))
    \qquad
    \text{as }\sens(f)\to\infty.
\]
In the more explicit form stated by Scheder and Tan \cite{ScTa13}, the conjecture predicts absolute constants \(K>0\) and \(\delta<1\) such that \(\AS(f)\le K\cdot\sens(f)^\delta\) for every monotone \(f\). A natural way to track progress is through the exponent
\[
    \alpha(f_n)
    =
    \liminf_{n\to\infty}
    \frac{\log\AS(f_n)}{\log\sens(f_n)}.
\]
The Servedio--Tan question asks whether \(\alpha(f_n)<1\) for every monotone sequence with \(\sens(f_n)\to\infty\). In the opposite direction, O'Donnell and Servedio \cite{OS07} gave a recursive construction achieving
\[
    \AS(f_n)
    \ge
    \sens(f_n)^{\alpha_{\mathrm{OS}}-o(1)},
    \qquad
    \alpha_{\mathrm{OS}}=\log_2(6-2\sqrt 5)\approx 0.6115,
\]
by composing a four-variable base gadget. As noted in \cite{ScTa13}, this was the previously recorded lower bound on \(\alpha\) for monotone functions in the literature.

\begin{remark}
As a side observation of this note, a more modest improvement over the O'Donnell--Servedio exponent is also obtained from the Fano-plane line-containment function. This monotone function has sensitivity \(3\) and average sensitivity \(63/32\), so recursive composition gives exponent
\[
    \frac{\log(63/32)}{\log 3}\approx 0.6166,
\]
slightly larger than \(\alpha_{\mathrm{OS}}\approx 0.6115\). We do not develop this finite gadget further: the main construction below already gives the qualitatively stronger exponent \(2/3\).
\end{remark}

Our main result is the following.

\begin{theorem}\label{thm:main}
There is a sequence of monotone Boolean functions \(H_k\) with \(\sens(H_k)\to\infty\) such that
\[
    \liminf_{k\to\infty}\frac{\log\AS(H_k)}{\log\sens(H_k)}\ge \frac{2}{3}.
\]
\end{theorem}

The argument has three layers. For each prime \(Q\) we build a monotone DNF \(R_{p,Q}\) over the variable set \(\F_Q\times\F_Q\) whose terms are graphs of polynomials of degree less than \(d\) for a suitable \(d\le Q\). Sharp first- and second-moment bounds on the number of satisfied terms give, via Paley--Zygmund, a \(p\)-biased influence of order \(Q\) on a fixed interval of biases. Combining two such blocks --- an outer block \(T\) at the bias induced by the complementary block, composed with the dual of an inner block \(B\) at bias \(1/2\) --- yields a monotone function with \(\Inf\ge cQ^2\) and \(\sens\le Q^3\). Letting \(Q\) range over an unbounded sequence of primes, the resulting sequence satisfies \(\liminf\log\AS/\log\sens\ge 2/3\).

\subsection*{Acknowledgements}
The arguments in this paper were developed in a harness combining two large language models: ChatGPT~5.5~Pro was used to generate and refine the mathematical arguments, and Claude~Opus~4.7 was used for the organization and formalization of the writeup. The author thanks @Liam06972452 on Twitter for catching that the extra millisecond of improvement claimed in an earlier draft of this note was not justified; the exponent here is correspondingly improved to \(2/3\approx 0.667\) instead of \(0.668\).

\section{Definitions and elementary composition facts}\label{sec:defs}

All Boolean functions are written in \(0/1\) notation. If \(V\) is a finite set, we identify an input \(x\in\{0,1\}^V\) with the set
\[
    X=\{v\in V:x_v=1\}.
\]
A function \(f:\{0,1\}^V\to\{0,1\}\) is monotone if \(X\subseteq Y\) implies \(f(X)\le f(Y)\).

For \(p\in(0,1)\), let \(X\sim\mu_p\) denote the random subset of \(V\) in which each element is included independently with probability \(p\). The \(p\)-biased total influence is
\[
    \Inf_p(f)
    =
    \sum_{v\in V}
    \Prb_{X\sim\mu_p}\bigl[f(X)\ne f(X\triangle\{v\})\bigr].
\]
The uniform total influence is the average sensitivity:
\[
    \AS(f)=\Inf_{1/2}(f).
\]
For an input \(X\), let \(\sens(f,X)\) be the number of coordinates whose flip changes the value of \(f\), and put
\[
    \sens(f)=\max_X \sens(f,X).
\]
For monotone functions it is useful to separate the two sides:
\[
    \sens_0(f)=
    \max_{f(X)=0}
    \bigl|\{v\notin X:f(X\cup\{v\})=1\}\bigr|,
\]
\[
    \sens_1(f)=
    \max_{f(X)=1}
    \bigl|\{v\in X:f(X\setminus\{v\})=0\}\bigr|.
\]
Then
\[
    \sens(f)=\max\{\sens_0(f),\sens_1(f)\}.
\]

The monotone dual of \(f\) is
\[
    f^\dagger(X)=1-f(V\setminus X).
\]
Then
\[
    \sens_0(f^\dagger)=\sens_1(f),
    \qquad
    \sens_1(f^\dagger)=\sens_0(f),
\]
and
\[
    \Inf_p(f^\dagger)=\Inf_{1-p}(f).
\]

If \(f:\{0,1\}^r\to\{0,1\}\) and \(g:\{0,1\}^m\to\{0,1\}\), their block composition \(f\circ g\) is obtained by replacing each input bit of \(f\) by an independent copy of \(g\). If
\[
    q=\Prb[g(X)=1]
\]
under the uniform measure, then
\begin{equation}\label{eq:inf-composition}
    \Inf(f\circ g)=\Inf_q(f)\Inf(g).
\end{equation}
Indeed, a variable inside the \(i\)-th copy of \(g\) is pivotal exactly when it is pivotal for that copy of \(g\), and the \(i\)-th coordinate of the outer function is pivotal at the vector of block outputs. The other block outputs are independent \(q\)-biased bits, and summing over all variables gives \eqref{eq:inf-composition}.

For monotone functions,
\begin{equation}\label{eq:sens-composition}
    \sens_0(f\circ g)\le \sens_0(f)\sens_0(g),
    \qquad
    \sens_1(f\circ g)\le \sens_1(f)\sens_1(g).
\end{equation}
For example, at a zero-input of \(f\circ g\), only upward-pivotal coordinates of the outer function can matter, and each corresponding block contributes at most \(\sens_0(g)\) upward-pivotal inner variables. The one-side case is identical.

\section{The polynomial-graph DNF block}\label{sec:block}

Fix a prime \(Q\), integers \(L\le Q\) and \(d\le L\), and a subset \(E\subseteq\F_Q\) with \(|E|=L\). The variable set is
\[
    V=E\times\F_Q.
\]
For every polynomial \(P\in\F_Q[z]\) with \(\deg P<d\), its graph over \(E\) is
\[
    \Gamma(P)=\bigl\{(t,P(t)):t\in E\bigr\}\subseteq V,
\]
so \(|\Gamma(P)|=L\). The polynomial-graph DNF is
\[
    R(X)=1
    \quad\Longleftrightarrow\quad
    \Gamma(P)\subseteq X
    \text{ for some }P\in\F_Q[z],\ \deg P<d.
\]
This is a monotone DNF whose terms all have size \(L\).

\begin{proposition}[Algebraic block]\label{prop:block}
Fix constants \(0<a<b<1\). There are constants \(c_0,c_1>0\) and a prime threshold \(Q_0\), depending only on \(a,b\), such that for every prime \(Q\ge Q_0\) and every \(p\in[a,b]\), one may choose \(L\le Q\), \(d<L\), and \(E\) so that the resulting block \(R_{p,Q}\) satisfies
\[
    c_0\le \Prb_{\mu_p}[R_{p,Q}=1]\le \frac14,
    \qquad
    \Inf_p(R_{p,Q})\ge c_1 Q,
\]
\[
    \sens_1(R_{p,Q})\le Q,
    \qquad
    \sens_0(R_{p,Q})\le Q^2.
\]
\end{proposition}

\begin{proof}
Write \(A=\log(1/p)\) and \(D_a=\exp(a^{-1}-1)\). Choose an integer \(K\) so large that
\begin{equation}\label{eq:K-choice}
    D_a\, b^{K-1}\le \frac14,
    \qquad
    b^{K-1}\le \frac14.
\end{equation}
For sufficiently large \(Q\), define
\[
    d=\Bigl\lfloor \tfrac{(Q-K-2)A}{\log Q}\Bigr\rfloor,
    \qquad
    L=\Bigl\lfloor \tfrac{d\log Q}{A}+K\Bigr\rfloor.
\]
Then \(Q/2\le L\le Q\) and \(0<d\le L-1\) for all sufficiently large \(Q\). Choose any \(E\subseteq\F_Q\) with \(|E|=L\); the DNF \(R_{p,Q}\) is defined above.

\emph{Sensitivity bounds.} Every term of the DNF has size \(L\le Q\), so deleting a variable outside a satisfied term cannot destroy that witness; hence
\[
    \sens_1(R_{p,Q})\le L\le Q.
\]
The total number of variables is \(|V|=LQ\le Q^2\), so trivially
\[
    \sens_0(R_{p,Q})\le |V|\le Q^2.
\]

\emph{Density bounds.} Let \(Z(X)\) be the number of satisfied polynomial-graph terms. There are \(Q^d\) polynomials of degree less than \(d\), each with a graph of size \(L\), so
\[
    \lambda:=\E_{\mu_p} Z=Q^d p^L.
\]
By the choice of \(L\), \(L-(d\log Q)/A\in[K-1,K]\), so
\begin{equation}\label{eq:lambda-bounds}
    a^K\le \lambda\le b^{K-1}.
\end{equation}
For distinct polynomials \(P,R\) of degree less than \(d\), the graphs intersect in \(|\Gamma(P)\cap\Gamma(R)|=|\{t\in E:(P-R)(t)=0\}|\) points, the root set of \(S:=P-R\) inside \(E\); since \(\deg S<d\) and \(S\ne 0\), this set has size at most \(d-1\). Substituting and reindexing \(R\) as \(P+S\),
\[
    \E_{\mu_p}[Z(Z-1)]
    =
    \sum_{P}\sum_{R\ne P}p^{|\Gamma(P)\cup\Gamma(R)|}
    =
    Q^d\, p^{2L}\!\!\sum_{\substack{S\ne 0\\ \deg S<d}} p^{-z(S)},
\]
where \(z(S):=|\{t\in E:S(t)=0\}|\). Setting \(\theta=p^{-1}-1\) and expanding
\[
    p^{-z(S)}=(1+\theta)^{z(S)}=\sum_{T\subseteq\{t\in E:S(t)=0\}}\!\!\theta^{|T|},
\]
the inner sum is bounded by
\[
    \sum_{\substack{S\ne 0\\ \deg S<d}} p^{-z(S)}
    \le
    \sum_{j=0}^{d-1}\binom{L}{j}\theta^{j}\, Q^{d-j},
\]
since for each \(T\subseteq E\) with \(|T|=j<d\) the nonzero polynomials of degree less than \(d\) vanishing on \(T\) form a set of size at most \(Q^{d-j}\), and any such \(S\) has at most \(d-1\) roots so \(j\le d-1\). Bounding the binomial sum by extending it,
\[
    \sum_{j=0}^{d-1}\binom{L}{j}\theta^{j} Q^{d-j}
    \le
    Q^d\bigl(1+\tfrac{\theta}{Q}\bigr)^{L}
    \le
    Q^d\exp(\theta L/Q)
    \le
    Q^d\, D_a,
\]
using \(L\le Q\) and \(\theta\le a^{-1}-1\). Therefore
\[
    \E_{\mu_p}[Z(Z-1)]\le D_a\,\lambda^2.
\]
Combining with~\eqref{eq:lambda-bounds} and~\eqref{eq:K-choice}, \(\E[Z(Z-1)]\le \tfrac14\lambda\). Since \(Z\) is integer-valued,
\[
    \Prb_{\mu_p}[Z=1]\ge \E Z-\E[Z(Z-1)]\ge \tfrac34\lambda\ge \tfrac34 a^K,
\]
and Markov's inequality gives
\[
    \Prb_{\mu_p}[R_{p,Q}=1]=\Prb_{\mu_p}[Z\ge 1]\le \E Z=\lambda\le b^{K-1}\le \tfrac14.
\]
Set \(c_0:=\tfrac34 a^K\).

\emph{Influence.} On the event \(Z=1\) there is a unique satisfied graph \(\Gamma(P)\); each of its \(L\) variables is pivotal at \(X\). Therefore
\[
    \Inf_p(R_{p,Q})\ge L\,\Prb_{\mu_p}[Z=1]\ge \tfrac{Q}{2}\cdot \tfrac34 a^K\ge c_1\,Q
\]
with \(c_1:=\tfrac38 a^K\).
\end{proof}

\section{Two-block composition}\label{sec:composition}

Apply Proposition~\ref{prop:block} on an interval containing \(1/2\), for instance \([1/2,2/3]\). Let \(B_Q:=R_{1/2,Q}\). Then, for absolute constants \(\alpha,c,C>0\),
\begin{equation}\label{eq:first-block}
    \alpha\le \Prb[B_Q=1]\le \tfrac14,
    \qquad
    \Inf(B_Q)\ge cQ,
\end{equation}
\begin{equation}\label{eq:first-block-sens}
    \sens_1(B_Q)\le Q,
    \qquad
    \sens_0(B_Q)\le Q^2.
\end{equation}
After shrinking \(\alpha\) if necessary, assume \(\alpha<1/4\). Set
\[
    G_Q:=B_Q^\dagger,
    \qquad
    \theta_Q:=\Prb[G_Q=1]=1-\Prb[B_Q=1]\in[\tfrac34,1-\alpha].
\]
By the duality identities,
\[
    \Inf(G_Q)=\Inf(B_Q)\ge cQ,
    \qquad
    \sens_0(G_Q)\le Q,
    \qquad
    \sens_1(G_Q)\le Q^2.
\]

Now apply Proposition~\ref{prop:block} again on the fixed interval \([\tfrac34,1-\alpha]\) (which lies inside \((0,1)\)). For each sufficiently large prime \(Q\), let \(T_Q:=R_{\theta_Q,Q}\). Then \(\Inf_{\theta_Q}(T_Q)\ge c'Q\), \(\sens_1(T_Q)\le Q\), and \(\sens_0(T_Q)\le Q^2\), for an absolute constant \(c'>0\).

Define the composed block
\[
    H_Q:=T_Q\circ G_Q,
\]
which is monotone. Since the output bias of \(G_Q\) is \(\theta_Q\), the influence-composition identity~\eqref{eq:inf-composition} gives
\begin{equation}\label{eq:H-influence}
    \Inf(H_Q)=\Inf_{\theta_Q}(T_Q)\,\Inf(G_Q)\ge c''Q^2,
\end{equation}
where \(c'':=c'c>0\) is an absolute constant. The sensitivity-composition bounds~\eqref{eq:sens-composition} give
\[
    \sens_0(H_Q)\le \sens_0(T_Q)\sens_0(G_Q)\le Q^2\cdot Q=Q^3,
\]
\[
    \sens_1(H_Q)\le \sens_1(T_Q)\sens_1(G_Q)\le Q\cdot Q^2=Q^3,
\]
so
\begin{equation}\label{eq:H-sens}
    \sens(H_Q)\le Q^3.
\end{equation}

\section{The exponent \(2/3\)}\label{sec:limit}

We now read off the asymptotic exponent.

\begin{proof}[Proof of Theorem~\ref{thm:main}]
Let \((Q_k)_{k\ge 1}\) be a sequence of primes with \(Q_k\to\infty\), and let \(H_k:=H_{Q_k}\) be the composed block constructed in Section~\ref{sec:composition}. By~\eqref{eq:H-influence} and~\eqref{eq:H-sens},
\[
    \AS(H_k)\ge c''\,Q_k^{2},
    \qquad
    \sens(H_k)\le Q_k^{3}.
\]
The universal bound \(\AS(f)\le \sens(f)\) (the pointwise sensitivity is at most its maximum) gives \(\sens(H_k)\ge \AS(H_k)\ge c''Q_k^2\), so \(\sens(H_k)\to\infty\). For \(k\) large enough that \(\sens(H_k)>1\),
\[
    \frac{\log\AS(H_k)}{\log\sens(H_k)}
    \ge
    \frac{2\log Q_k+\log c''}{3\log Q_k}
    =
    \frac{2}{3}+\frac{\log c''}{3\log Q_k}.
\]
Taking \(\liminf\) as \(k\to\infty\),
\[
    \liminf_{k\to\infty}\frac{\log\AS(H_k)}{\log\sens(H_k)}\ge \frac{2}{3}.
\]
\end{proof}

\begin{thebibliography}{ScTa13}

\bibitem[FHH$^+$14]{FHH+14}
Y.~Filmus, H.~Hatami, S.~Heilman, E.~Mossel, R.~O'Donnell, S.~Sachdeva, A.~Wan, and K.~Wimmer,
\emph{Real analysis in computer science: A collection of open problems},
Simons Institute for the Theory of Computing, 2014.

\bibitem[OS07]{OS07}
R.~O'Donnell and R.~A.~Servedio,
Learning monotone decision trees in polynomial time,
\emph{SIAM J.~Comput.}\ 37(3):827--844, 2007.

\bibitem[ScTa13]{ScTa13}
D.~Scheder and L.-Y.~Tan,
On the average sensitivity and density of \(k\)-CNF formulas,
in \emph{Approximation, Randomization, and Combinatorial Optimization (APPROX/RANDOM 2013)},
Lecture Notes in Computer Science, vol.\ 8096, Springer, Berlin, Heidelberg, 2013, pp.\ 683--698.

\end{thebibliography}

\end{document}
